% TOP_PROBLEM: {"problemId": "problem.infinitude-of-mersenne-primes", "canonicalTitle": "Infinitude of Mersenne Primes", "releaseRank": 261, "primaryDomain": "number_theory_arithmetic_geometry", "primaryDomainLabel": "Number theory & arithmetic geometry", "displayStatus": "open", "releaseStatus": "open"}
% !TeX program = lualatex
% ATTEMPT_STATUS: unresolved
\documentclass[11pt]{article}
\usepackage[margin=25mm]{geometry}
\usepackage{fontspec}
\setmainfont{Latin Modern Roman}
\usepackage{amsmath,amssymb,mathtools}
\usepackage{unicode-math}
\setmathfont{Latin Modern Math}
\usepackage{xurl,hyperref}
\hypersetup{colorlinks=true,urlcolor=blue}
\setcounter{secnumdepth}{0}
\setlength{\parindent}{0pt}
\setlength{\parskip}{5pt}
\setlength{\emergencystretch}{3em}
\begin{document}
\section*{\texorpdfstring{Infinitude of Mersenne Primes}{Infinitude of Mersenne Primes}}\label{261}
\subsection{Definitions and mathematical statement}
A Mersenne number is $M_p=2^p-1$. The selected infinitude question is
\[
 \forall B>0\quad\exists p>B\quad
 p\text{ prime and }2^p-1\text{ prime}.
\]
Primality of the exponent is necessary, but not sufficient, for a Mersenne number to be prime. The catalogue's equivalent even-perfect-number formulation uses the Euclid--Euler correspondence
$N=2^{p-1}(2^p-1)$ between Mersenne primes and even perfect numbers. The question requires neither a density prediction nor an efficient procedure guaranteed to find the next example. Arbitrarily large verified examples alone do not prove infinitude without a general argument.
\par\smallskip\textit{Formulation and concept references:} \href{https://t5k.org/mersenne/index.html}{[S1]}.

\subsection{Short English statement}
Are there infinitely many primes of the form $2^p-1$ with $p$ prime, equivalently infinitely many even perfect numbers?
\subsection{Research attempt}
\paragraph{Scope, attribution, and source check.}
This attempt was developed by GPT-6 Astra Ultra on 27 September 2026.
The source [S1] distinguishes prime exponents from prime Mersenne
numbers and explains the connection with even perfect numbers.
The primary computational project [E1] describes the Lucas--Lehmer
test; an independently checkable implementation is supplied with this
attempt. No claim about the current largest example is needed.
The work below proves the test, proves infinitude of distinct prime
divisors across prime exponents, and isolates why neither statement
proves infinitude of Mersenne primes themselves.

\paragraph{Exponent reduction and pairwise coprimality.}
If $n=ab$ with $a,b>1$, then
\[
 2^n-1=(2^a-1)(1+2^a+\cdots+2^{a(b-1)}),                \tag{1}
\]
so a prime Mersenne number must have prime exponent. Conversely,
$2^{11}-1=23\cdot89$, showing that the exponent condition is not
sufficient.

The Euclidean algorithm on exponents gives
\[
              \gcd(2^a-1,2^b-1)=2^{\gcd(a,b)}-1.        \tag{2}
\]
To see this, when $a\ge b$ subtract $2^{a-b}(2^b-1)$ from
$2^a-1$ to replace the first entry by $2^{a-b}-1$, and iterate.
The terminal zero exponent gives the claimed gcd. Hence $M_p$ and
$M_q$ are coprime for different primes $p,q$.
Each is greater than one, so choosing a prime divisor from each yields
infinitely many distinct ordinary primes. This is a complete infinitude
argument for divisors, but its conclusion does not require even one
of those $M_p$ to be prime.

\paragraph{Restrictions on every prime factor.}
Let $p$ be an odd prime and $q\mid2^p-1$ a prime. The multiplicative
order of $2$ modulo $q$ divides $p$ and cannot be one, so it equals $p$.
Thus $p\mid q-1$. Since $q$ is odd and $p$ is odd, actually
\[
                              q=2kp+1.                  \tag{3}
\]
Moreover $2^{(q-1)/2}=(2^p)^k\equiv1\pmod q$.
Euler's criterion and the supplementary law for $2$ give
\[
                              q\equiv\pm1\pmod8.        \tag{4}
\]
Equations (3)--(4) greatly reduce trial division, but the allowable
factors still extend all the way to $\sqrt{M_p}$ for detecting
compositeness. Their absence below a fixed cutoff is insufficient.

For example $p=11$ allows $q=23=2p+1\equiv-1\pmod8$.
More generally, if $p\equiv3\pmod4$ and $q=2p+1$ is prime, then
$q\equiv7\pmod8$ and Euler's criterion implies
$2^p\equiv1\pmod q$. For $p>3$, this is a proper factor of $M_p$.
This gives a conditional family of composite prime-exponent Mersenne
numbers; it does not establish infinitude of such exponents because
the required primality pattern for $2p+1$ has not been proved infinite.

\paragraph{An exact primality certificate.}
For an odd prime $p$, define integers
\[
                 s_0=4,\qquad s_{j+1}=s_j^2-2.
\]
We prove
\[
            M_p\text{ is prime}\quad\Longleftrightarrow\quad
                           s_{p-2}\equiv0\pmod{M_p}.    \tag{5}
\]
The exponent $p=2$ is handled separately, with $M_2=3$.
The recurrence may be performed modulo $M_p$ at every step.

First suppose the right side of (5). Let $q$ be any prime divisor of
$M_p$. It is odd and is not three, since $2^p-1\equiv1\pmod3$.
In a field containing a square root $t$ of $3$ over $\mathbb F_q$, put
$\alpha=2+t$. Then $(2+t)(2-t)=1$, and induction gives
\[
                      s_j=\alpha^{2^j}+\alpha^{-2^j}.    \tag{6}
\]
The zero residue at $j=p-2$ implies
$\alpha^{2^{p-1}}=-1$, and hence $\alpha$ has order exactly $2^p$.
If $3$ is a square modulo $q$, this order divides $q-1$.
If not, the quadratic-field Frobenius sends $t$ to $-t$, so
$\alpha^q=\alpha^{-1}$ and the order divides $q+1$.
In either case $q+1\ge2^p$, hence $q\ge M_p$.
Since $q\mid M_p$, equality follows: $M_p=q$ is prime.
This sufficiency argument does not assume primality of the modulus.

Conversely suppose $q=M_p$ is prime. Since $p\ge3$ is odd,
$q\equiv7\pmod{12}$ and $q\equiv7\pmod8$.
The quadratic-character laws give $(3/q)=-1$ and $(2/q)=1$.
Work in $\mathbb F_{q^2}$ with $t^2=3$, so $t^q=-t$.
Let $\beta=1+t$. Its norm is
$\beta^{q+1}=(1+t)(1-t)=-2$, and
$\beta^2=2(2+t)=2\alpha$. Therefore
\[
 \alpha^{(q+1)/2}
   =\frac{\beta^{q+1}}{2^{(q+1)/2}}
   =\frac{-2}{2}=-1.                                   \tag{7}
\]
Here the denominator is $2$ because Euler's criterion gives
$2^{(q-1)/2}=1$. Since $(q+1)/2=2^{p-1}$, writing
$z=\alpha^{2^{p-2}}$ gives $z^2=-1$, and so
$z+z^{-1}=0$. Equation (6) proves the necessary residue in (5).
All divisions used here are legitimate because $q$ is odd.

The standard finite-field Frobenius and quadratic-character laws are
the algebraic inputs; the order argument and recurrence are written
out to prevent the certificate from becoming a black-box numerical
primality claim. This is the classical Lucas--Lehmer mechanism,
not a newly discovered test.

\paragraph{Even perfect numbers, including the converse.}
If $M_p$ is prime, multiplicativity of the divisor sum gives
\[
 \sigma(2^{p-1}M_p)=(2^p-1)(M_p+1)=2(2^{p-1}M_p).
\]
Thus the displayed number is perfect.
Conversely let an even perfect number be $N=2^{p-1}m$ with $m$ odd
and $p\ge2$ an integer. The equality $\sigma(N)=2N$ yields
\[
                    (2^p-1)\sigma(m)=2^p m.
\]
Coprimality forces $m=(2^p-1)a$ and $\sigma(m)=2^p a=m+a$.
Both $m$ and $a$ are divisors of $m$. If $a>1$, then $1,a,m$
are distinct divisors and their sum exceeds $m+a$, a contradiction.
Thus $a=1$ and $\sigma(m)=m+1$, which means $m$ is prime.
Equation (1) then forces $p$ prime. This establishes the precise
equivalence in the question without assumptions on odd perfect numbers.

\paragraph{Why finite sieves and Euclid's argument stop short.}
For each fixed odd prime $q$, the condition $q\mid M_p$ at a prime
exponent $p$ requires $p=\operatorname{ord}_q(2)$. Thus a fixed finite
list of trial factors excludes at most finitely many prime exponents.
Passing beyond those exponents only shows that all factors are new;
their sizes are unbounded and may still be smaller than $M_p$.
Pairwise coprimality (2) makes this issue particularly transparent.

Likewise multiplying previously found Mersenne primes and adding one
produces an integer with a new prime factor, but supplies no reason
that the integer or that factor is of the form $2^p-1$.
The special exponential shape is not preserved by Euclid's construction.
Probability models predicting a divergent expected count do not
establish infinitely many zero residues in (5); the moduli and the
lengths of the recurrences both depend on $p$.

\paragraph{Verification and the exact quantified gap.}
The accompanying exact-integer check runs (5) for all prime exponents
$p\le127$, recording the successful exponents and the nonzero residues
otherwise. At small exponents it also compares with direct trial
division. The proof of (5), rather than the number of tested cases,
is what certifies each successful modulus. The script additionally
tests (2), the factor restrictions, and the first even perfect numbers.
All computational conclusions are confined to their stated ranges.

A route to infinitude using (5) must prove the genuinely uniform statement
\[
 \forall B\ \exists\text{ prime }p>B:
          s_{p-2}\equiv0\pmod{2^p-1}.
\]
The certificate decides each instance but provides no lower bound for
the count of successful instances. This distinction cannot be removed
by making the computation longer.

\paragraph{Final status.}
\textbf{UNRESOLVED.} The complete partial results are the primality
criterion, the Euclid--Euler correspondence, and infinitude of distinct
prime factors among the $M_p$. No infinitude proof for prime values is
obtained. The first gap is a lower bound tending to infinity for the
number of successful prime exponents. Concrete next targets are a
rigorous distribution statement for the Lucas--Lehmer residues, control
of all potential factors up to $\sqrt{M_p}$ for infinitely many $p$,
or a structure theorem forcing prime rather than composite values.
None is supplied by the finite-factor sieve developed here.

\subsection{Sources}
\begin{itemize}
\item[S1]"Mersenne Primes: History, Theorems and Lists," The PrimePages (t5k.org), maintained curated reference; printed from the PrimePages, (c) Reginald McLean.
\url{https://t5k.org/mersenne/index.html}.
\textit{Formulation source.}
\item[E1]Great Internet Mersenne Prime Search, \emph{The Math}, project documentation.
\url{https://www.mersenne.org/various/math.php}.
\textit{Primary computational context for Mersenne testing.}
\item[E2]D. H. Lehmer, \emph{On Lucas's Test for the Primality of Mersenne's Numbers}, Journal of the London Mathematical Society s1-10 (1935), 162--165.
\url{https://doi.org/10.1112/jlms/s1-10.2.162}.
\textit{Primary historical reference for the test proved above.}
\end{itemize}
\end{document}
